\documentclass{article}
\usepackage{amsmath,amssymb}
\usepackage{xurl}
\usepackage[hidelinks]{hyperref}
\setlength{\emergencystretch}{2em}
% Shared commands for the notes in docs/paper-gaps/.
% Every note loads this file via % Shared commands for the notes in docs/paper-gaps/.
% Every note loads this file via % Shared commands for the notes in docs/paper-gaps/.
% Every note loads this file via \input{command}. Project-wide notation
% belongs here, never in per-note redefinitions.

% --- Project configuration ------------------------------------------------
\newcommand{\projectIssueBase}{https://github.com/LionSR/QICLean/issues/}
\newcommand{\projectPrBase}{https://github.com/LionSR/QICLean/pull/}
\newcommand{\projectDocBase}{https://sirui-lu.com/QICLean/docs/}
% ---------------------------------------------------------------------------

\DeclareUnicodeCharacter{2080}{\ifmmode{}_{0}\else\textsubscript{0}\fi}
\DeclareUnicodeCharacter{2081}{\ifmmode{}_{1}\else\textsubscript{1}\fi}
\DeclareUnicodeCharacter{2082}{\ifmmode{}_{2}\else\textsubscript{2}\fi}
\DeclareUnicodeCharacter{2099}{\ifmmode{}_{n}\else\textsubscript{n}\fi}

\newcommand{\Lean}{\textsc{Lean}}

% A formal declaration name, upright sans, breakable at underscores.
\ExplSyntaxOn
\NewDocumentCommand{\leanid}{m}{
  \ifmmode
    \textsf{\detokenize{#1}}
  \else
    \group_begin:
    \urlstyle{sf}
    \tl_set:Nn \l_tmpa_tl {#1}
    \tl_replace_all:Nnn \l_tmpa_tl {\_} {_}
    \exp_args:NV \path \l_tmpa_tl
    \group_end:
  \fi
}
\ExplSyntaxOff

% Traceability links; use in footnotes, not in the running prose.
\newcommand{\ghissue}[1]{\href{\projectIssueBase#1}{GitHub issue \##1}}
\newcommand{\ghpr}[1]{\href{\projectPrBase#1}{PR \##1}}
\newcommand{\doclink}[1]{\href{\projectDocBase#1.html}{\path{#1}}}

% Domain notation, using \providecommand so a note-local refinement stays
% possible.
\providecommand{\tr}{\operatorname{tr}}

% Dirac notation and the identity operator, as in blueprint/src/macros.
\usepackage{dsfont}
\providecommand{\ket}[1]{|#1\rangle}
\providecommand{\bra}[1]{\langle #1|}
\providecommand{\braket}[2]{\langle #1 | #2 \rangle}
\providecommand{\ketbra}[2]{|#1\rangle\!\langle #2|}
\providecommand{\Id}{\mathds{1}}
\providecommand{\one}{\mathds{1}}
\providecommand{\C}{\mathbb{C}}
\providecommand{\MN}[1]{M_{#1}(\C)}
\providecommand{\MD}{M_D(\C)}

% Verdict marker: \gapnote{<kind>}{<status>}. Kind is the policy
% classification (clarification, local-correction, scope-restriction,
% unfaithful, false-source, open-gap); status is open, wip (elimination
% underway), resolved, or historical. Machine-read by the site index;
% prints a verdict line.
\newcommand{\gapnote}[2]{\par\noindent\textbf{Verdict:} #1 (#2).\par}
. Project-wide notation
% belongs here, never in per-note redefinitions.

% --- Project configuration ------------------------------------------------
\newcommand{\projectIssueBase}{https://github.com/LionSR/QICLean/issues/}
\newcommand{\projectPrBase}{https://github.com/LionSR/QICLean/pull/}
\newcommand{\projectDocBase}{https://sirui-lu.com/QICLean/docs/}
% ---------------------------------------------------------------------------

\DeclareUnicodeCharacter{2080}{\ifmmode{}_{0}\else\textsubscript{0}\fi}
\DeclareUnicodeCharacter{2081}{\ifmmode{}_{1}\else\textsubscript{1}\fi}
\DeclareUnicodeCharacter{2082}{\ifmmode{}_{2}\else\textsubscript{2}\fi}
\DeclareUnicodeCharacter{2099}{\ifmmode{}_{n}\else\textsubscript{n}\fi}

\newcommand{\Lean}{\textsc{Lean}}

% A formal declaration name, upright sans, breakable at underscores.
\ExplSyntaxOn
\NewDocumentCommand{\leanid}{m}{
  \ifmmode
    \textsf{\detokenize{#1}}
  \else
    \group_begin:
    \urlstyle{sf}
    \tl_set:Nn \l_tmpa_tl {#1}
    \tl_replace_all:Nnn \l_tmpa_tl {\_} {_}
    \exp_args:NV \path \l_tmpa_tl
    \group_end:
  \fi
}
\ExplSyntaxOff

% Traceability links; use in footnotes, not in the running prose.
\newcommand{\ghissue}[1]{\href{\projectIssueBase#1}{GitHub issue \##1}}
\newcommand{\ghpr}[1]{\href{\projectPrBase#1}{PR \##1}}
\newcommand{\doclink}[1]{\href{\projectDocBase#1.html}{\path{#1}}}

% Domain notation, using \providecommand so a note-local refinement stays
% possible.
\providecommand{\tr}{\operatorname{tr}}

% Dirac notation and the identity operator, as in blueprint/src/macros.
\usepackage{dsfont}
\providecommand{\ket}[1]{|#1\rangle}
\providecommand{\bra}[1]{\langle #1|}
\providecommand{\braket}[2]{\langle #1 | #2 \rangle}
\providecommand{\ketbra}[2]{|#1\rangle\!\langle #2|}
\providecommand{\Id}{\mathds{1}}
\providecommand{\one}{\mathds{1}}
\providecommand{\C}{\mathbb{C}}
\providecommand{\MN}[1]{M_{#1}(\C)}
\providecommand{\MD}{M_D(\C)}

% Verdict marker: \gapnote{<kind>}{<status>}. Kind is the policy
% classification (clarification, local-correction, scope-restriction,
% unfaithful, false-source, open-gap); status is open, wip (elimination
% underway), resolved, or historical. Machine-read by the site index;
% prints a verdict line.
\newcommand{\gapnote}[2]{\par\noindent\textbf{Verdict:} #1 (#2).\par}
. Project-wide notation
% belongs here, never in per-note redefinitions.

% --- Project configuration ------------------------------------------------
\newcommand{\projectIssueBase}{https://github.com/LionSR/QICLean/issues/}
\newcommand{\projectPrBase}{https://github.com/LionSR/QICLean/pull/}
\newcommand{\projectDocBase}{https://sirui-lu.com/QICLean/docs/}
% ---------------------------------------------------------------------------

\DeclareUnicodeCharacter{2080}{\ifmmode{}_{0}\else\textsubscript{0}\fi}
\DeclareUnicodeCharacter{2081}{\ifmmode{}_{1}\else\textsubscript{1}\fi}
\DeclareUnicodeCharacter{2082}{\ifmmode{}_{2}\else\textsubscript{2}\fi}
\DeclareUnicodeCharacter{2099}{\ifmmode{}_{n}\else\textsubscript{n}\fi}

\newcommand{\Lean}{\textsc{Lean}}

% A formal declaration name, upright sans, breakable at underscores.
\ExplSyntaxOn
\NewDocumentCommand{\leanid}{m}{
  \ifmmode
    \textsf{\detokenize{#1}}
  \else
    \group_begin:
    \urlstyle{sf}
    \tl_set:Nn \l_tmpa_tl {#1}
    \tl_replace_all:Nnn \l_tmpa_tl {\_} {_}
    \exp_args:NV \path \l_tmpa_tl
    \group_end:
  \fi
}
\ExplSyntaxOff

% Traceability links; use in footnotes, not in the running prose.
\newcommand{\ghissue}[1]{\href{\projectIssueBase#1}{GitHub issue \##1}}
\newcommand{\ghpr}[1]{\href{\projectPrBase#1}{PR \##1}}
\newcommand{\doclink}[1]{\href{\projectDocBase#1.html}{\path{#1}}}

% Domain notation, using \providecommand so a note-local refinement stays
% possible.
\providecommand{\tr}{\operatorname{tr}}

% Dirac notation and the identity operator, as in blueprint/src/macros.
\usepackage{dsfont}
\providecommand{\ket}[1]{|#1\rangle}
\providecommand{\bra}[1]{\langle #1|}
\providecommand{\braket}[2]{\langle #1 | #2 \rangle}
\providecommand{\ketbra}[2]{|#1\rangle\!\langle #2|}
\providecommand{\Id}{\mathds{1}}
\providecommand{\one}{\mathds{1}}
\providecommand{\C}{\mathbb{C}}
\providecommand{\MN}[1]{M_{#1}(\C)}
\providecommand{\MD}{M_D(\C)}

% Verdict marker: \gapnote{<kind>}{<status>}. Kind is the policy
% classification (clarification, local-correction, scope-restriction,
% unfaithful, false-source, open-gap); status is open, wip (elimination
% underway), resolved, or historical. Machine-read by the site index;
% prints a verdict line.
\newcommand{\gapnote}[2]{\par\noindent\textbf{Verdict:} #1 (#2).\par}

\title{Spectral windows and the cutoff in hard truncation}
\author{}
\date{2026-10-29}
\begin{document}
\maketitle
\gapnote{false-source}{wip}

\section{The conditional spectral-window comparison}

The second item of Lemma \texttt{lem:hard-trunc} of
Arad--Landau--Vazirani--Vidick,
\href{https://arxiv.org/abs/1602.08828}{arXiv:1602.08828v2}, source lines
926--948, states that the low-energy subspace $\tilde T_{\eta'}$ of the
truncated Hamiltonian $\tilde H_t$ is $\delta$-close to the low-energy
subspace $T_\eta$ of $H$. Its proof (lines 953--969) bounds the mass above
$\varepsilon_0+\eta'$ of a vector supported on eigenvalues of $H$ in a
closed window $[\lambda-h,\lambda+h]$ with $\lambda\le\varepsilon_0+\eta$
(lines 955--962), splits a vector $\psi\in T_\eta$ into such pieces, and sums
the bounds with the Cauchy--Schwarz inequality. Given valid restricted
perturbation estimates, this spectral-window argument can be completed with
the corrections below. The unrestricted physical truncated-side estimate
also needs a cutoff qualification, as shown in the next section.

\emph{The number of windows.} Line 963 bounds the number of windows by
$\lceil(\eta-\varepsilon_0)/(2h)\rceil$, which subtracts an energy from a
width. The intended count is $\lceil\eta/(2h)\rceil$: closed windows with
centres $\varepsilon_0+\eta-h-2hk$, $0\le k<\lceil\eta/(2h)\rceil$, cover
$[\varepsilon_0,\varepsilon_0+\eta]$, and every centre is at most
$\varepsilon_0+\eta$, as line 955 requires. Line 967 then bounds the
Cauchy--Schwarz factor by $\sqrt{\eta/(2h)}$, where the number of windows
gives $\sqrt{\lceil\eta/(2h)\rceil}$.

\emph{The ground energy of the truncation.} The proof bounds the mass above
$\varepsilon_0+\eta'$, whereas $\tilde T_{\eta'}$ is defined from the ground
energy $\tilde\varepsilon_0\le\varepsilon_0$ of $\tilde H_t$. Since larger
subspaces are closer, it suffices to enlarge $\eta'$ by
$\varepsilon_0-\tilde\varepsilon_0$, which the first item bounds by
$Ce^{-c(t-\eta)}$.

\emph{Two meanings of $\delta$.} Line 962 uses $\delta$ for the distance in
Markov's inequality, and the lemma uses it for the closeness parameter.

\emph{The sign of $\xi$.} Equation \texttt{eq:def-xi} prints $\xi$ with a
positive exponent in $t-\eta$; the intended quantity is
$\xi=e^{-\Omega(t-\eta)}$, the bound of Theorem~2.6 of Arad--Kuwahara--Landau,
\href{https://arxiv.org/abs/1406.3898}{arXiv:1406.3898}, Theorem
\texttt{thm:lowspec}.

The formal statement \leanid{Submodule.isCloseTo_spectralSpanLE_of_perturbation}
is the closeness step. It is stated for two self-adjoint operators diagonal
in given orthonormal bases, an energy $c$ in place of $\varepsilon_0+\eta$, a
lower bound $m<c$ of the eigenvalues of $H$ below $c$ in place of
$\varepsilon_0$, and a threshold $c'\ge c$. With
$\|(H-\tilde H_t)\varphi\|\le\xi\|\varphi\|$ on $T(c)$ as a hypothesis, it
concludes that $\tilde T(c')$ is $\delta$-close to $T(c)$ when
$\lceil(c-m)/(2h)\rceil(\xi+h)^2\le\delta(c'-c+h)^2$. The extra $h$ comes from
centring the top window at $c-h$. For $h=\xi$ this holds when
$c'-c\ge2\xi\sqrt{\lceil(c-m)/(2\xi)\rceil/\delta}-\xi$, which is at most
$\sqrt{2\xi(c-m)/\delta}+2\xi/\sqrt\delta$. Adding the shift
$\varepsilon_0-\tilde\varepsilon_0$ gives a valid $\eta'$, which has the form
$\eta+\sqrt{\eta/\delta}\,e^{-\Omega(t-\eta)}$ of the lemma for $\eta\ge1$ and
$\delta\le1$.

\section{A flat sector inside the requested low-energy space}

The sign of $\xi$ is not the only qualification needed for the physical
estimate. Proposition~4.4(1) of
\href{https://arxiv.org/pdf/1602.08828v2}{ALVV, arXiv:1602.08828v2},
printed page~15, bounds
\begin{align*}
    \|(H-\widetilde H_t)\widetilde\Pi_{\le\varepsilon_0+\eta}\|
     & \le\xi.
\end{align*}
The preceding paragraph of Section~4.2 says that $t$ is chosen sufficiently
large relative to $\eta$, but the proposition gives no numerical cutoff
condition and Lemma~4.5 explicitly states $t>\eta$.
Definition~2.5 and Theorem~2.6 of
\href{https://arxiv.org/pdf/1406.3898}{AKL, arXiv:1406.3898v3},
printed pages~8--9, likewise have no explicit condition keeping the
requested spectral space below the clipped flat sector.

Here is an actual nearest-neighbour chain for which $t>\eta$ alone is
insufficient. Let $K\ge2$, take qubits on sites $0,\ldots,K+1$, and write
$P_b(i)$ for the on-site projection onto label $b$.
For $1\le i<K$, set $h_i=P_1(i)$. The remaining bonds are
\begin{align*}
    h_0 & =(I-P_1(0)/4)P_0(1),        &
    h_K & =\tfrac12(P_1(K)+P_1(K+1)).
\end{align*}
There is exactly one positive contraction on each bond $\{i,i+1\}$.
Both the crossing term $h_0$ and the last internal term depend on both
of their sites. For $J=\{1,\ldots,K+1\}$, put
\begin{align*}
    A            & =H_J=\sum_{i=1}^K h_i,   &
    H            & =A+h_0,                  &
    \widetilde H & =\min(A,\tfrac32 I)+h_0.
\end{align*}
The ground energy of $A$ is zero. The vector
$\Omega=|1,0,\ldots,0\rangle$ attains ground energy $3/4$ for both
$H$ and $\widetilde H$. Indeed, a configuration with label zero at
site~$1$ has boundary energy at least $3/4$; a configuration with label
one there has internal energy at least one, both before and after clipping.

For the unit vector $\psi=|1,\ldots,1\rangle$,
\begin{align*}
    H\psi                      & =K\psi,        &
    \widetilde H\psi           & =\tfrac32\psi, &
    \widetilde\Pi_{\le7/4}\psi & =\psi.
\end{align*}
Thus, with the fixed parameters $t=3/2>\eta=1$,
\begin{align*}
    \|(H-\widetilde H)\widetilde\Pi_{\le\varepsilon_0+\eta}\|
     & \ge K-\tfrac32.
\end{align*}
Every length-independent constant is exceeded by an Archimedean choice
of $K$. This applies to both the printed growing expression for $\xi$ and
the intended decaying expression, since their values are fixed here.
Only the truncated-side bound is refuted by this example; the
original-side projection is not used.

The corresponding actual-chain proofs are source-written in TNLean's
\path{TNLean/MPS/Examples/HardTruncationCutoff.lean}, with hosted acceptance
pending. The physical application with a sufficiently large explicit cutoff
remains a separate result. The generic conditional spectral-window theorem
above is unchanged and does not assert the physical estimate without its
restricted perturbation premise. This boundary is tracked in
\href{https://github.com/LionSR/TNLean/issues/9844}{TNLean issue 9844}.
\end{document}
